\chapter{Core, Crop Engine, and Viewer}
\docsource{docs/design/02-viewer-interaction.md}
\docsource{docs/design/03-core-architecture.md}

\section{Core Boundary}

\texttt{core}는 geometry, crop semantics, point source abstraction, spatial
index를 소유한다. viewer와 app은 core API를 호출하지만, core는 windowing,
OpenGL, CLI parser, 포맷 codec을 알지 않는다.

\section{Crop Semantics}

모든 box는 OBB로 처리한다. membership test는 point를 box-local frame으로
변환한 뒤 slab/AABB test를 수행한다.

\begin{lstlisting}[language=C++]
local = inverse(box.rotation) * (point - box.center);
inside = abs(local.x) <= half.x &&
         abs(local.y) <= half.y &&
         abs(local.z) <= half.z;
\end{lstlisting}

multi-box 조합은 include box의 union/intersection을 먼저 계산하고, exclude box는
항상 subtract한다. 이 규칙은 viewer preview와 exported crop이 동일해야 하는
load-bearing contract다.

\section{Spatial Index Policy}

\begin{lstlisting}[language=C++]
enum class IndexMode { Never, Always, Auto };

struct IndexPolicy {
  IndexMode mode = IndexMode::Auto;
  size_t auto_point_threshold = 200'000;
  size_t expected_queries = 1;
  bool interactive = false;
};
\end{lstlisting}

one-shot CLI crop은 brute-force가 더 단순하고 빠를 수 있다. interactive session,
반복 query, 큰 cloud에서는 octree를 만들어 box range candidate를 줄인다.
threshold는 benchmark로 보정한다.

\section{Viewer Contract}

viewer는 full-resolution cloud를 소유하지 않는다. core가 full-res 데이터를
유지하고, viewer는 budgeted decimated copy를 렌더링한다. LOD는 display-only다.
실제 crop은 항상 full-res core 데이터에 대해 실행한다.

\begin{longtable}{p{0.25\linewidth}p{0.65\linewidth}}
\toprule
\textbf{Concern} & \textbf{결정} \\
\midrule
Renderer & OpenGL 3.3 core profile, GLFW, Dear ImGui docking, ImGuizmo. \\
Camera & orbit/pan/zoom 중심. point picking은 KD-tree나 screen-space candidate로
처리한다. \\
Box editing & ImGuizmo transform handle과 bounds handle을 함께 제공한다. \\
Responsiveness & load, index, crop, write는 worker thread에서 실행하고 UI는
progress/cancel/result queue만 처리한다. \\
\bottomrule
\end{longtable}

\section{Concurrency}

UI thread는 blocking 작업을 수행하지 않는다. pipeline stage는
\texttt{common::ThreadPool}, \texttt{Progress}, \texttt{stop\_token}을 통해
진행률과 취소를 보고한다. CLI는 같은 stage를 동기적으로 await하고, GUI는 queue로
결과를 drain한다.

